% This file is part of the modular source for Stratified Relational Workspace Theory.
% !TeX root = ../Stratified_Relational_Workspace_Theory.tex
\theorypart{XV}{Boundary Formation and Stability}

\topic{50}{Candidate Workspace Boundaries}

A boundary is not merely a subset of physical components. It is a subset
relative to a declared node decomposition, prediction target, intervention
family, environment distribution, and time scale.

\begin{definition}[Boundary problem]\label{def:boundary-problem}
A boundary problem is a tuple
\[
\mathfrak B
=
\left(
V,
\mathfrak S,
\tau,
\mathcal Y,
\mathcal I,
\Pi,
\mathcal G_B
\right),
\]
where \(\mathfrak S\subseteq2^V\) is the admissible candidate family,
\(\tau\) is the temporal scale, \(\mathcal Y\) is the operational target,
\(\mathcal I\) is the feasible intervention family, \(\Pi\) is the
environment distribution, and \(\mathcal G_B\) is a symmetry group under
which boundaries are considered equivalent.
\end{definition}

Unless the application requires otherwise, the candidate family excludes
\[
\varnothing,
\qquad
V,
\qquad
\{i\}\ \text{for }i\in V.
\]
This anti-triviality rule prevents the empty set, the entire realization, or a
single centralized controller from winning solely through a degenerate score.

\subtopic{50.1 Observational recurrence, dependence, and outward influence}

For \(S\in\mathfrak S\), define
\[
\Ret_\tau(S)
=
I\!\left(
X_S(t+\tau);X_S(t)
\mid X_{\bar S}(t),E_t
\right),
\]
\[
\Dep_\tau(S)
=
I\!\left(
X_S(t+\tau);X_{\bar S}(t)
\mid X_S(t),E_t
\right),
\]
\[
\Out_\tau(S)
=
I\!\left(
X_{\bar S}(t+\tau);X_S(t)
\mid X_{\bar S}(t),E_t
\right).
\]
These quantities are observational. They can be nonzero because of common
causes, selection, or feedback paths and therefore cannot alone establish a
causal boundary.

\subtopic{50.2 Causal boundary terms}

Let \(\iota_{S,x}\) be a feasible intervention that changes the internal
workspace coordinate of \(S\) to \(x\) while matching background variables and
intervention cost. Define outward causal efficacy
\[
\boxed{
\COut_\tau(S)
=
\mathbb E_{x,x'\sim\nu_S}
D\!\left(
\mathcal L^{\iota_{S,x}}
(X_{\bar S}(t+\tau)),
\mathcal L^{\iota_{S,x'}}
(X_{\bar S}(t+\tau))
\right).
}
\]
Define inward causal dependence analogously by intervening on \(\bar S\) and
measuring the future law of \(X_S\):
\[
\boxed{
\CIn_\tau(S)
=
\mathbb E_{u,u'\sim\nu_{\bar S}}
D\!\left(
\mathcal L^{\iota_{\bar S,u}}
(X_S(t+\tau)),
\mathcal L^{\iota_{\bar S,u'}}
(X_S(t+\tau))
\right).
}
\]
Neither term should be minimized to zero in general. A useful workspace may
receive substantial external input. The boundary objective must instead
express the intended balance between internal recurrence, external dependence,
and outward efficacy.

\subtopic{50.3 Structural and anti-centralization gates}

For the induced subsystem on \(S\), define
\[
\mathsf{Coh}(S)
=
\operatorname{tr}
\left(
P_{S,d}L_SP_{S,d}
\right),
\]
where \(P_{S,d}\) projects onto the selected low-energy workspace modes. Let
\[
\mathsf{Cut}(S)
=
\sum_{e:s(e)\in S,\ t(e)\notin S}w_e
+
\sum_{e:s(e)\notin S,\ t(e)\in S}w_e.
\]
Let \(k_\varepsilon(S)\) be the minimal sufficient coalition size inside the
candidate subsystem. A candidate is admissible only if
\[
\operatorname{SRW}_{\Theta_S}(W_S)=1,
\qquad
k_\varepsilon(S)\geq2,
\]
unless the claim explicitly concerns a centralized workspace.

\topic{51}{Boundary Objective, Recovery, and Persistence}

A general score is
\[
\boxed{
\begin{aligned}
\mathcal B_t(S)
={}&
 a_R\Ret_\tau(S)
+a_O\Out_\tau(S)
+a_C\COut_\tau(S)\\
&-a_D\Dep_\tau(S)
-a_I\CIn_\tau(S)
-a_L\mathsf{Coh}(S)\\
&-a_K\mathsf{Cut}(S)
-a_Z\mathsf{Size}(S)
-a_T\Instability_t(S).
\end{aligned}
}
\]
The coefficients and all normalizations are part of the boundary problem. No
single choice is universal. The selected boundary is
\[
\widehat S_t
\in
\argmax_{S\in\mathfrak S_{\mathrm{pass}}}
\widehat{\mathcal B}_t(S),
\]
where \(\mathfrak S_{\mathrm{pass}}\) contains only candidates that satisfy the
required workspace and anti-centralization gates.

\subtopic{51.1 Boundary equivalence and identifiability}

If a symmetry \(g\in\mathcal G_B\) leaves the full observational and
interventional law invariant, \(S\) and \(gS\) cannot be distinguished by the
declared experiment. The identifiable boundary is then the orbit
\[
[S]_{\mathcal G_B}
=
\{gS:g\in\mathcal G_B\}.
\]
A unique subset should not be claimed when only an equivalence class is
identifiable.

\subtopic{51.2 Stable boundary recovery \(T_{17}\) }

\begin{theorem}[Held-out boundary recovery]\label{thm:heldout-boundary-recovery}
Let \(\mathfrak S_{\mathrm{pass}}\) be finite and let
\[
S_t^*
=
\argmax_{S\in\mathfrak S_{\mathrm{pass}}}
\mathcal B_t(S)
\]
be unique. Define the population margin
\[
\Delta_t^B
=
\mathcal B_t(S_t^*)
-
\max_{S\neq S_t^*}\mathcal B_t(S)>0.
\]
If
\[
\sup_{S\in\mathfrak S_{\mathrm{pass}}}
\left|
\widehat{\mathcal B}_t(S)-\mathcal B_t(S)
\right|
\leq\varepsilon_t,
\]
then every empirical maximizer \(\widehat S_t\) satisfies
\[
\boxed{
\mathcal B_t(S_t^*)-
\mathcal B_t(\widehat S_t)
\leq2\varepsilon_t.
}
\]
If \(2\varepsilon_t<\Delta_t^B\), then
\[
\boxed{\widehat S_t=S_t^*.}
\]
\end{theorem}

\begin{proof}
Apply Theorem~\ref{thm:uniform-structure-selection} to
\(J(S)=-\mathcal B_t(S)\).
\end{proof}

The theorem converts uncertainty in mutual-information estimation, causal
interventions, structural scores, and stability penalties into an explicit
boundary-identification condition. It also exposes when the data are
insufficient: a small margin makes the boundary intrinsically fragile.

\subtopic{51.3 Temporal boundary stability}

\begin{theorem}[Boundary persistence under score drift]
\label{thm:boundary-drift-stability}
Assume \(S_t^*\) is the unique maximizer of \(\mathcal B_t\) with margin
\(\Delta_t^B\), and suppose
\[
\sup_{S\in\mathfrak S_{\mathrm{pass}}}
|\mathcal B_{t+1}(S)-\mathcal B_t(S)|
\leq\delta_t.
\]
If
\[
2\delta_t<\Delta_t^B,
\]
then
\[
\boxed{S_{t+1}^*=S_t^*.}
\]
More generally,
\[
\mathcal B_{t+1}(S_{t+1}^*)
-
\mathcal B_{t+1}(S_t^*)
\leq2\delta_t.
\]
\end{theorem}

\begin{proof}
For every competitor \(S\neq S_t^*\),
\[
\begin{aligned}
\mathcal B_{t+1}(S_t^*)-
\mathcal B_{t+1}(S)
&\geq
\mathcal B_t(S_t^*)-
\mathcal B_t(S)-2\delta_t\\
&\geq
\Delta_t^B-2\delta_t>0.
\end{aligned}
\]
Hence the maximizer is unchanged. The approximate bound follows by comparing
both time-indexed scores at their respective maximizers.
\end{proof}

If the score depends on a parameter \(\vartheta_t\) and satisfies
\[
\sup_S
|\mathcal B_{\vartheta}(S)-\mathcal B_{\vartheta'}(S)|
\leq
L_Bd(\vartheta,\vartheta'),
\]
then the sufficient stability condition becomes
\[
2L_Bd(\vartheta_t,\vartheta_{t+1})
<
\Delta_t^B.
\]

\subtopic{51.4 Analyst-discovered, system-instantiated, and self-maintained boundaries}

Three claims must remain separate.

\begin{description}[leftmargin=2.1em,style=nextline]
\item[Analyst-discovered boundary.]
An external procedure fits \(\widehat S\) after observing the system. The
system need not represent or use that boundary.

\item[System-instantiated boundary.]
The system contains a channel, mask, or routing variable causally equivalent to
\(\widehat S\), and interventions on that variable alter downstream operation.

\item[Self-maintained boundary.]
The system detects degradation of its boundary certificate and takes actions
that restore membership, transfer, or reconstruction properties.
\end{description}

Only the second and third claims concern a boundary realized within the
system. A post hoc partition that predicts behavior well is not by itself a
self-discovered boundary.

\subtopic{51.5 Boundary formation certificate}

A boundary \(\widehat S\) is certified as endogenously formed only if:
\begin{enumerate}
\item the membership mask was not directly encoded in \(\Sigma_B\) or
 \(\Sigma_Y\);
\item \(e_B^{\mathrm{MDL}}\) and \(e_B^{\mathrm{causal}}\) exceed declared
 thresholds;
\item held-out recovery satisfies Theorem~\ref{thm:heldout-boundary-recovery};
\item the boundary remains stable under admissible environment drift or its
 changes are predicted by score-margin closure;
\item matched history-shuffle and label-leakage controls fail to reproduce
 the result, as do full-system, empty-set, singleton, and centralized
 boundary nulls;
\item the induced subsystem satisfies its workspace certificate; and
\item a system-instantiation intervention demonstrates that the boundary is
 used by the system rather than only by the analyst.
\end{enumerate}

The results below establish the statistical and deterministic stability components. They
does not prove that an arbitrary adaptive system will spontaneously learn such
a boundary.
